\documentclass[a4paper,11pt,dvipdfmx]{jsarticle}
\usepackage{amsmath, amssymb, amsthm}
\usepackage{bm}
\usepackage{mathtools}
\usepackage{comment}
\usepackage{algorithm}
\usepackage{algorithmic}
\usepackage{bm}
\usepackage{mathtools}
\usepackage{enumitem}
\mathtoolsset{showonlyrefs}
\usepackage[dvipdfmx, colorlinks=true, linkcolor=blue, citecolor=blue, urlcolor=blue, setpagesize=false]{hyperref}
\usepackage{pxjahyper}
% --- 定理・定義環境 ---
\theoremstyle{definition}
\newtheorem{definition}{定義}[section]
\newtheorem{theorem}[definition]{定理}
\newtheorem{lemma}[definition]{補題}
\newtheorem{assumption}{仮定}
\newtheorem{proposition}[definition]{命題}
\newtheorem{requirement}[definition]{要請}
\newtheorem{example}[definition]{例}
\newcommand{\fU}{\underline{f}}
\newcommand{\gU}{\underline{g}}
\newcommand{\uU}{\underline{u}}
\newcommand{\aU}{\underline{a}}
\newcommand{\bU}{\underline{b}}
\newcommand{\cU}{\underline{c}}
\newcommand{\dU}{\underline{d}}
\newcommand{\vU}{\underline{v}}
\newcommand{\zeroU}{\underline{0}}
\newcommand{\C}[2]{C_{{#1},{#2}}}
\newcommand{\cy}[3]{\mathcal{C}_{{#1},{#2},{#3}}}
\newcommand{\inv}{^{-1}}
\newcommand{\HHX}{\hat{H}_X}
\newcommand{\HHZ}{\hat{H}_Z}
\newcommand{\tr}{^\top}

% =========================================
\begin{document}

\title{ガースを大きくする研究}
\author{}
\date{\today}
\maketitle
\tableofcontents
\clearpage
\section{行列をより疎にする}
長さ4のサイクルを含まないようなプロトグラフを設計したい。
アクティブ行列の列重みは3,行重みは12とする。
このために必要なアクティブ行列のサイズを考える。
つまり、条件は
\begin{itemize}
  \item 列重み3
  \item 行重み12
  \item 各行の内積が1以下
\end{itemize}
今まで、行列の左半分と右半分は同じ構造を持っていたが、そうすると必ず長さ4のサイクルが含まれてしまう。
一旦行重み6で考え左右でずらすことでうまく長さ4のサイクルを避けられないか考える。
\subsection{プロトグラフの最小サイズの導出}
最小行数を$M$、最小列数を$N$とする。
列重みは3なので、各列が$\binom{3}{2} = 3$通りの行ペアの重なりを消費する。
全ての列が異なる行ペアの重なりを持つ必要があるので、
\begin{align}
  3N \leq \binom{M}{2} = \frac{M(M-1)}{2}
\end{align}
が成り立つ必要がある。
また、$12M=3N$より、$N=4M$である。
これらの条件を満たす最小の$M$は25である。
よって、最小のアクティブ行列のサイズは$25 \times 100$である。

\section{空間結合LDPC符号の構造}

空間結合LDPC符号のパリティ検査行列 $H_{\text{SC}}$ は、以下のような帯行列（Banded Matrix）構造を持つ。
ここで、$L$ は結合の長さ（Coupling Length）、$w$ は結合幅を表す。

\begin{equation}
H_{\text{SC}} = 
\begin{pmatrix}
H_0(0) &        &        &        & \\
H_1(0) & H_0(1) &        &        & \\
\vdots & H_1(1) & \ddots &        & \\
H_w(0) & \vdots & \ddots & H_0(L-1) & \\
       & H_w(1) &        & H_1(L-1) & \\
       &        & \ddots & \vdots   & \\
       &        &        & H_w(L-1) & 
\end{pmatrix}
\end{equation}

\bibliography{refs}
\bibliographystyle{unsrt} %参考文献出力スタイル
\end{document}